\chapter{ActStats}
\label{chap:actstats}

\section{Overview}

\actsim depends on the companion package \texttt{actstats}, authored by 
Juntao Zhang and Tingting Shi, which provides actuarial-style wrappers 
around standard statistical distributions. The \texttt{actstats} package 
exposes a familiar interface for fitting, sampling, and evaluating 
distributions, and it is the engine behind \actsim's distribution 
fitting and stochastic simulation classes.

The \texttt{actstats} package is imported under the alias
\texttt{act} throughout the manual:

\begin{lstlisting}[style=pythonstyle]
from actstats import actuarial as act
\end{lstlisting}

\section{Why a Dedicated Actuarial Wrapper?}

Many actuarial computations could be performed directly using
\texttt{scipy.stats} or \texttt{numpy.random}. However, parameter conventions 
in \texttt{scipy} and \texttt{numpy.random} do not always match standard 
actuarial textbook conventions. For example, \texttt{scipy.stats.lognorm} 
is parametrised by a shape parameter $s$ and optional location and scale 
arguments, whereas actuaries typically work with the underlying normal mean 
$\mu$ and standard deviation $\sigma$.

\texttt{actstats} provides actuarial-friendly parametrisations and
exposes a consistent API across distributions. It also provides
specialised utilities such as the non-homogeneous Poisson process used
for incurred date simulation in \actsim.

\section{Supported Distributions}

The distributions exposed through \texttt{actstats} that are used
internally by \actsim include:

\begin{itemize}[leftmargin=*]
    \item Continuous: \texttt{normal}, \texttt{logistic},
          \texttt{exponential}, \texttt{gamma}, \texttt{beta},
          \texttt{lognormal}, \texttt{weibull}, \texttt{pareto},
          \texttt{uniform}.
    \item Discrete: \texttt{poisson}, \texttt{negative\_binomial}.
\end{itemize}

These names match the entries used in the default \texttt{config.yaml}
shipped with \actsim and may be referenced from the configuration loader.

\section{Parameter Conventions}
The parameter conventions are described in Section~\ref{app:parameter_conventions}

\section{Common Operations}

The interface exposed by \texttt{actstats} mirrors familiar statistical
patterns:

\begin{lstlisting}[style=pythonstyle]
from actstats import actuarial as act

# Construct a frozen distribution
dist = act.lognormal(0.5, 0.2)

# Draw random samples
samples = dist.rvs(size=10000)

# Evaluate the probability density function
pdf_values = dist.pdf([1.0, 1.5, 2.0])

# Evaluate the cumulative distribution function
cdf_values = dist.cdf([1.0, 1.5, 2.0])

# Evaluate the inverse CDF (quantile function)
quantiles = dist.ppf([0.5, 0.9, 0.99])
\end{lstlisting}

\section{Discrete Distributions}

Frequency simulations rely on discrete distributions such as Poisson
and negative binomial. Sampling is performed in the same way:

\begin{lstlisting}[style=pythonstyle]
from actstats import actuarial as act

# Poisson with mean 10
freq_dist = act.poisson(10)
counts = freq_dist.rvs(size=1000)

# 80th percentile of Poisson(10)
q80 = act.poisson.ppf(0.8, 10)
\end{lstlisting}

\section{Non-Homogeneous Poisson Process}

\texttt{actstats} also provides a utility for simulating event times from
a non-homogeneous Poisson process (NHPP) with a sinusoidal seasonal
intensity:

\begin{lstlisting}[style=pythonstyle]
from actstats import actuarial as act

# Baseline intensity 10, seasonality amplitude 0.5,
# phase shift 0, exposure 1 year
nhpp = act.nonhomogeneous_poisson(10, 0.5, 0, 1)

# Generate event time fractions in [0, 1]
times = nhpp.rvs(n_events=20)
\end{lstlisting}

This utility is used internally by \texttt{ClaimSimulator} to assign
incurred dates to simulated claims (see
Chapter~\ref{chap:claim_simulator}).

\section{Date Conversion Helpers}

\texttt{actstats} also provides date conversion helpers, such as
\texttt{fraction\_to\_date\_full}, that map fractions of an exposure
period to calendar dates. These are used by \actsim when assigning
incurred and development dates to simulated claims.
